\section{Cuidados Éticos}
\label{sec:ethics}

This experience report respects the Brazilian regulatory framework for research involving human participants. The study is part of an umbrella project at the host university's Research Ethics Committee (Comit\^{e} de \'{E}tica em Pesquisa, CEP), approved under CAAE (omitted for double-blind review), and follows the requirements of CNS Resolutions 466/2012 \cite{cns466_2012}, 510/2016 \cite{cns510_2016}, and 674/2022 \cite{cns674_2022}, as well as Operational Norm 001/2013 \cite{cns_no001_2013}.

\subsection{Participant consent and rights}

Participation was voluntary, non-remunerated, and could be withdrawn at any moment without consequence. Before answering the questionnaire, each participant read an Informed Consent Form (\emph{Termo de Consentimento Livre e Esclarecido}, TCLE) approved by the ethics committee. The form described the purpose of the research, the activities involved, the foreseen risks and benefits, the right to withdraw at any time, the team's contact information, and the way the collected data would be handled. Consent had to be explicitly given before any further question could be answered.

The risks were considered minimal. The activity is comparable to using a website and answering a few open-ended reflection questions, and could cause at most some discomfort from the time spent or fatigue from reading. No clinical, physical, or invasive procedures were involved. The foreseen benefits are indirect: participants contributed to research on the teaching of Nielsen's heuristics and to the refinement of an educational artefact for the HCI community.

\subsection{Anonymity}

No personally identifying information was collected at any point. The demographic questions are all categorical (role, age range, area, prior familiarity with the heuristics, and practical experience with HCI evaluation), and no names, emails, IP addresses, or other identifiers are recorded by the game or by the questionnaire. In the analysis and in this paper, participants are referred to only as $P_1, P_2, \ldots$, with a short demographic tag added when context for a quote is needed. No participant images are used in this paper.

\subsection{Data protection (LGPD)}

The processing of participant data follows the Brazilian General Data Protection Law \cite{lgpd2018}. Because the questionnaire collects only categorical and self-reported responses with no direct identifiers, the resulting dataset is pseudonymous by design. Responses are stored in a spreadsheet linked to the form, with access restricted to the research team. Aggregate analyses and anonymised excerpts are used for academic communication; the raw data are not redistributed. The data are retained for five years from collection, in line with the consent form, and are used only for the purposes described to participants.
